In this assignment, we compare two neural models for news classification: a Convolutional Neural Network (CNN) and a Long Short-Term Memory network (LSTM). To make the comparison fair, both models use the same data splits, preprocessing pipeline, and evaluation setup. The overall classification task pipelines is illustrated in Figure~\ref{fig:pipeline}.



\paragraph{Shared input representation.}
Both models receive the same preprocessed and tokenized text (see Section~\ref{sec:dataset} for preprocessing and tokenization details). Sequences are padded or truncated to the chosen maximum length. Both models begin with a trainable embedding layer that maps token ids to dense vectors. Note that the CNN has inherent feature extraction through its convolutional layers, while the LSTM relies on the embeddings and recurrent connections for representation learning.

% \paragraph{CNN model.}
% The CNN is based on the text classification model of \cite{kim-2014-convolutional}. After the embedding layer, the model applies several 1-dimensional convolution filters in parallel. These filters are used to detect informative local patterns in the text, such as short phrases or combinations of nearby words. We use multiple kernel sizes so the model can capture patterns of different lengths. After the convolution step, a ReLU activation function is applied, followed by max-pooling to keep the strongest signal from each filter. The resulting features are then concatenated, followed by dropout for regularisation, and finally passed to a linear layer that predicts one of the four classes.

\paragraph{LSTM model.}
The LSTM model is based on the Long Short-Term Memory architecture introduced by \cite{hochreiter1997lstm}. It reads the input sequence token by token, which makes it suitable for modelling word order and longer dependencies in the text. We use a stacked 2-layer LSTM. To avoid learning from padding tokens, we handle variable-length sequences by skipping over padded positions during the recurrent computation, so the model only processes the real tokens in each example. For classification, we use the final hidden state of the top LSTM layer as the representation of the whole article. In the bidirectional setting, the model combines information from both directions before making the final prediction. Dropout is applied before the final linear layer. This model architecture was chosen as it was the best performing model in Assignment 2; however, due to the change in tokenization for a fair comparison, we re-tuned the dropout rate and learning rate to ensure a fair comparison with the BERT model. As training the LSTM is not the focus of this assignment, no further details on the training process will be provided, but hyperparameters used for final LSTM training can be seen in Table \ref{tab:hyperparameters}.

\paragraph{BERT model.}
The BERT model is based on the pre-trained Bidirectional Encoder Representations from Transformers (BERT) architecture \cite{devlin2019bertpretrainingdeepbidirectional}. We use the $bert-base-uncased$ variant, which has 12 layers, a hidden size of 768, and 12 attention heads, resulting in a total of 110 million parameters. The model is fine-tuned on our classification task by adding a feed-forward neural network classification head on top of the [CLS] token representation. The BERT model is trained end-to-end, meaning that the weights of the pre-trained BERT layers and the classification head are updated during training on our dataset. Before using the regular learning rate for fine-tuning, we apply a warmup phase of 500 steps where the learning rate gradually increases from zero to the specified learning rate. This helps stabilize training and prevents large updates to the pre-trained weights in the early stages of fine-tuning. As in LSTM, dropout is applied at the classification head to prevent overfitting.

\paragraph{Hyperparameters.}
Table~\ref{tab:hyperparameters} shows the hyperparameter search space for BERT model and the best hyperparameters for LSTM from Assignment 2. The best setting for BERT was selected based on macro-F1 on the development set.

% TODO: add the search spaces for BERT
\begin{table}[H]
\centering
\small
\begin{tabular}{lll}
\toprule
\textbf{Parameter} & \textbf{BERT} & \textbf{LSTM} \\
\midrule
Embedding dim $d$       & ---                & 256 \\
Hidden dim $H$          & 768                & 256 \\
Sequence length         & 128                & 128 \\
Learning rate           & 2e-5               & 3e-4 \\
Weight decay            & 0.01               & 0.01 \\
Dropout                 & ---                & 0.4 \\
Layers $L$              & 12                 & 2 \\
Bidirectional           & True               & True \\
Attention heads $A$     & 12                 & --- \\
Head dim                & 256                & --- \\
Head dropout            & 0.1                & --- \\
\bottomrule
\end{tabular}
\caption{Hyperparameter search space for the BERT and best hyperparameters for the LSTM model.}
\label{tab:hyperparameters}
\end{table}
% LSTM
% pooling_type=mean
% pack_sequences=false
% num_train_epochs=8
% metric_for_best_model=eval_loss (greater_is_better=false)
% eval_strategy=epoch, save_strategy=epoch

\paragraph{Shared training settings.}
Both models were trained with the Adam optimiser. We used gradient clipping with a maximum norm of 1.0 and early stopping with a patience of 2 epochs based on validation loss. The batch size for BERT was set to 16, while for the LSTM it was set to 32 to boost the training times.

% The sequence length was set to 128 for both models, and the final training runs lasted for at most 15 epochs. During hyperparameter tuning, each configuration was trained for 5 epochs.

\paragraph{Robustness/Slicing: Keyword Masking Probe.}
As the first robustness/slicing study for BERT model, we implemented a simple keyword masking probe to test how much the models rely on specific keywords for classification. We identified the top 15 most informative keywords for each class based on the training data using TF-IDF scores. Then, we created a modified test set where these keywords were masked out (replaced with a special [MASK] token). We evaluated the models on this masked test set to see how much their performance drops compared to the original test set. A significant drop in performance would suggest that the model relies heavily on these keywords for classification, while a smaller drop would indicate that the model is using more contextual information and is less sensitive to specific keywords. This analysis helps us understand the robustness of the models to changes in key lexical cues and whether they are learning more generalizable patterns or just memorizing specific terms.

\paragraph{Robustness/Slicing: Label-Noise Sensitivity (Simple).}
As the second robustness/slicing study for BERT model, we implemented a simple label-noise sensitivity analysis. We reduce the training set sizes to 25\%, 50\%, and 75\% of the original training set, and evaluate the models trained on these smaller subsets on the full test set. This allows us to see how sensitive the models are to the amount of training data, their learning efficiency, and generalization capabilities. If the model's performance drops significantly with less training data, it may indicate that the model is heavily reliant on large amounts of data to learn effectively. Conversely, if the performance remains relatively stable even with reduced training data, it suggests that the model is able to learn useful representations and generalize well even with limited data.

